\documentclass[10pt]{article}
\usepackage[margin=1in]{geometry}
\usepackage[T1]{fontenc}
\usepackage{lmodern,microtype,amsmath,amssymb,amsthm}
\theoremstyle{definition}
\newtheorem{assumption}{Assumption}
\begin{document}
\section*{Stationarity for Nonnegative Smooth Sample Losses}
We seek a stationary point of a possibly nonconvex objective $f:E\to\mathbb{R}$ on a finite-dimensional Euclidean space $E$. The objective is the expectation of nonnegative sample losses with $L$-Lipschitz gradients, where $L>0$, and $f(0)\leq\Delta$ for a given $\Delta>0$. Each oracle call returns the gradients of one fresh, independently sampled loss at both queried points. For a tolerance $\varepsilon>0$, the goal is a measurable output $\bar{x}\in E$ satisfying $\mathbb{E}[\|\nabla f(\bar{x})\|]\leq\varepsilon$, with expectation over the oracle samples and the algorithm's independent randomness.

The parameters satisfy $2^{32}\varepsilon\leq\sqrt{L\Delta}$. For a nonnegative integer $r$ specifying the extra width, set
\[
\begin{aligned}
n&=\left\lfloor\frac{\sqrt{L\Delta}}{2^{32}\varepsilon}\right\rfloor,\\
m&=[512(n+1)]^2,\\
d&=2^{21}(2n^3+1)n m^2+r+1.
\end{aligned}
\]
The optimization problem is
\[
\inf_{x\in E} f(x),
\]
where $E=(\mathbb{R}^{m\times d})^n$ carries the Euclidean norm $\|x\|^2=\sum_{j=1}^n\|x_j\|_{\mathrm F}^2$, with $\|\cdot\|_{\mathrm F}$ the Frobenius norm. The algorithm receives $E$, the sample space $\Omega=\{1,\ldots,m\}$, the starting point $x^0=0$, and the parameters $L,\Delta,\varepsilon$.

Our assumptions:
\begin{assumption}[Sample loss regularity]
For each $\omega\in\Omega$, the loss $\ell_\omega:E\to\mathbb{R}$ is differentiable. Writing $g(x,\omega)=\nabla\ell_\omega(x)$, we require for every $x,y\in E$ and $\omega\in\Omega$ that
\[
\begin{aligned}
\ell_\omega(x)&\geq 0,\\
\|g(x,\omega)-g(y,\omega)\|&\leq L\|x-y\|.
\end{aligned}
\]
The map $(x,\omega)\mapsto g(x,\omega)$ is jointly measurable.
\end{assumption}

\begin{assumption}[Population objective and initial value]
There is a probability distribution $\nu$ on $\Omega$ such that $\omega\mapsto\ell_\omega(x)$ is $\nu$-integrable for every $x\in E$, and
\[
\begin{aligned}
f(x)&=\mathbb{E}_{\omega\sim\nu}[\ell_\omega(x)] &&\text{for every }x\in E,\\
f(x^0)&\leq\Delta.
\end{aligned}
\]
\end{assumption}

On call $k$, the algorithm may query any pair $(x,y)\in E\times E$ and receives $(g(x,\xi_k),g(y,\xi_k))$. The samples $(\xi_k)_{k\geq 1}$ are independent random variables with common law $\nu$, and the sample sequence is independent of the algorithm's randomness. After the allotted calls, the output $\bar{x}$ must be measurable, its gradient norm must be integrable, and
\[
\mathbb{E}\bigl[\|\nabla f(\bar{x})\|\bigr]\leq\varepsilon.
\]
\end{document}
